% Independent proof supplied by /root/homotopy_levels, 2026-11-19.
% Replacement proof for C4.02; retain the existing section identifier.
For possibly overlapping dominated convex bands put
\(Y=X_J^*\), \(L_0=X_L^*\), and \(Z=Y\cup L_0\). Define the relative path space
\[
 P_{J,L}=\{\gamma\in C([0,1],Z):\gamma(0)\in Y,
                    \ \gamma(t)\in L_0\text{ for every }t>0\}.
\]
It contains \emph{all} such paths: paths starting in \(Y\cap L_0\) need not be constant. In particular it contains all constant paths on that intersection. Give it the compact-open topology. The set \(Y\) is a closed pure union of strata in the induced finite metric HSS \(Z\), so \cite[Proposition 3.2]{Quinn1988} supplies a fixed neighborhood \(N\) of \(Y\) and a nearly strict stratum-preserving deformation
\(h:N\times[0,1]\to Z\), with \(h_0=\mathrm{id}\), \(h_1=r:N\to Y\), and \(h_t|_Y=\mathrm{id}\). For \(t<1\), a point outside \(Y\) remains in its original stratum.

For each open neighborhood \(U\) of \(Y\) in \(Z\), let \(P_U\) be the paths in \(P_{J,L}\) whose whole image lies in \(U\). There is a natural zigzag
\[
 \operatorname{const} P_{J,L}
 \ \longleftarrow\ \{P_U\}_{U\supset Y}
 \ \xrightarrow{e_1}\ \{U\cap L_0\}_{U\supset Y}=V_J^L.
\]
The equality on the right means equality of pro-germs: restrictions of the unions of the finitely many individual stratum neighborhoods form a cofinal system of neighborhoods of \(Y\) in \(Z\). Indeed each such union is open and contains \(Y\); conversely an open neighborhood of \(Y\) in \(Z\) is the restriction of an ambient open set containing \(Y\), and intersecting that ambient open set with the prescribed open complements of the lower frontiers gives the required individual neighborhoods.

\textbf{The left arrow.} The path space is metrizable, using the uniform metric of any bounded compatible metric on \(Z\), and hence is paracompact. The open subsets
\[
 C_n=\{\gamma:\gamma([0,1/n])\subset U\},\qquad n\ge1,
\]
cover \(P_{J,L}\). Choose a locally finite partition of unity subordinate to a refinement of this cover; for each member of the refinement retain its integer \(n\). If its partition functions are \(\lambda_a\) and integers are \(n_a\), put
\(\epsilon(\gamma)=\sum_a\lambda_a(\gamma)/n_a\).
This is positive and continuous. Every active \(1/n_a\) is an admissible initial interval for \(\gamma\); since their finite maximum is also admissible, \(\gamma([0,\epsilon(\gamma)])\subset U\). Thus
\[
 S\gamma(u)=\gamma(\epsilon(\gamma)u)
\]
is a continuous map \(P_{J,L}\to P_U\). The homotopy
\[
 K_s\gamma(u)=\gamma\bigl((1-s+s\epsilon(\gamma))u\bigr)
\]
joins the identity to inclusion followed by \(S\), remains in \(P_{J,L}\), and restricts to a homotopy in \(P_U\) when \(\gamma\in P_U\). It fixes the initial point. Consequently every inclusion \(P_U\hookrightarrow P_{J,L}\) is a homotopy equivalence. These inclusions are a natural transformation, so the left arrow is a pro-homotopy equivalence. The endpoints are not fixed by this shortening homotopy; their coherent comparison in \(L_0\) is the explicitly displayed terminal segment of \(\gamma\).

\textbf{The right arrow.} For a target neighborhood \(U\), put
\[
 W_U=\{x\in N:h(x,[0,1])\subset U\}.
\]
Compactness of \([0,1]\) makes \(W_U\) an open neighborhood of \(Y\). Define
\[
 s_U:W_U\cap L_0\longrightarrow P_U,
 \qquad s_U(x)(u)=h(x,1-u).
\]
For \(x\in L_0\setminus Y\), the positive-time points of this reversed path remain in the stratum of \(x\), hence in \(L_0\). For \(x\in Y\cap L_0\) the path is constant. Thus it belongs to the relative path space in either case, and \(e_1s_U=\mathrm{id}\). All the maps \(s_U\) are restrictions of the same formula using the fixed \(h\); for finitely many targets a common smaller neighborhood gives literal agreement. They therefore define one pro map.

For the other composite, take \(\gamma\in P_{W_U}\) and define, for \(t>0\),
\[
 H_t(\gamma)(u)=
 \begin{cases}
 h\bigl(\gamma(t),1-u/t\bigr),&0\le u\le t,\\
 \gamma(u),&t\le u\le1,
 \end{cases}
 \qquad H_0(\gamma)=\gamma.
\]
The formulas agree at \(u=t\). Every positive-time point lies in \(L_0\), including when \(\gamma(t)\in Y\cap L_0\), because then \(h\) fixes that point. Every path lies in \(U\), and its endpoint is always \(\gamma(1)\). At \(t=1\) the path is \(s_U(e_1\gamma)\). Continuity away from \((t,u)=(0,0)\) is immediate. At that corner, \(\gamma(t)\to\gamma(0)\in Y\), and continuity of \(h\), compactness of its time interval, and \(h(\gamma(0),v)=\gamma(0)\) uniformly in \(v\) prove continuity. The same argument with \(\gamma\) varying gives continuity of the adjoint map on \(P_{W_U}\times[0,1]^2\), hence continuity into the compact-open path space. This proves the second inverse homotopy after shrinking. All refinements restrict this one fixed formula. The right arrow is consequently a pro-homotopy equivalence, in fact over \(L_0\).

The zigzag thus gives a canonical equivalence of pro-homotopy types. Its arrows literally commute with the endpoint maps to \(L_0\); the start map on the middle term is evaluation at zero. Under the inverse supplied by \(h\), the induced map to \(Y\) is \(r\), and the displayed \(H_t\) supplies its coherent comparison with the initial point of any representative path.

When \(J\cap L=\varnothing\), this relative path space is the conventional strict holink of \((Z,Y)\), recovering Quinn's usual formulation. When \(J=L\ne\varnothing\), that conventional strict holink is empty, whereas \(V_J^J=X_J^*\). The relative path space instead becomes \(C([0,1],X_J^*)\), which contracts by endpoint evaluation to the constant paths. Thus the literal strict-holink claim needs disjointness; the all-paths relative formulation above is valid with overlap.
